\documentclass[11pt]{article}

\usepackage[a4paper,margin=1in]{geometry}
\usepackage[T1]{fontenc}
\usepackage[utf8]{inputenc}
\usepackage{lmodern}
\usepackage{microtype}
\usepackage{amsmath,amssymb,amsthm,mathtools}
\usepackage{booktabs}
\usepackage{hyperref}
\usepackage[authoryear,round]{natbib}

\hypersetup{
  colorlinks=true,
  linkcolor=blue,
  citecolor=blue,
  urlcolor=blue
}

\title{\textbf{A Learning-Controlled Elephant Random Walk:}\\
Martingale Structure, Diffusive Limits and Reproducible Evidence}

\author{Prepared as a research manuscript}
\date{\today}

\newtheorem{theorem}{Theorem}
\newtheorem{proposition}{Proposition}
\newtheorem{corollary}{Corollary}
\newtheorem{lemma}{Lemma}
\newtheorem{assumption}{Assumption}
\theoremstyle{definition}
\newtheorem{definition}{Definition}
\theoremstyle{remark}
\newtheorem{remark}{Remark}

\newcommand{\F}{\mathcal{F}}
\newcommand{\E}{\mathbb{E}}
\newcommand{\Pbb}{\mathbb{P}}
\newcommand{\R}{\mathbb{R}}
\newcommand{\1}{\mathbf{1}}
\newcommand{\Var}{\operatorname{Var}}
\newcommand{\Law}{\mathcal{L}}

\begin{document}
\maketitle

\begin{abstract}
The elephant random walk is a non-Markovian stochastic process whose next
increment is drawn from its complete history and is then repeated or reversed.
This article introduces a learning-controlled elephant random walk (LC-ERW), in
which the memory parameter is no longer fixed but is selected online by a
predictable learning rule. The construction is intended as a mathematically
tractable bridge between long-memory random walks and artificial-intelligence
mechanisms for adaptive exploration. The main result is a martingale
normalisation valid for arbitrary predictable memory policies. Under a stable
learning condition, namely convergence of the learned memory coefficient to a
diffusive limit smaller than one half, the LC-ERW satisfies a central limit
theorem with the same variance form as an elephant random walk with the limiting
memory coefficient. The proof is based on triangular martingale convergence and
explicit control of the predictable quadratic variation. Reproducible R
experiments are provided to check the variance prediction and to illustrate the
transition from adaptive behaviour to the limiting Gaussian regime.
\end{abstract}

\noindent\textbf{Keywords:} elephant random walk; artificial intelligence;
online learning; martingales; central limit theorem; non-Markovian random walk.

\section{Introduction}

Random walks with memory provide simple models in which local moves are affected
by a long historical record. The elephant random walk (ERW), introduced by
\citet{schutz2004elephants}, is a canonical example: at each time the walker
selects one of its past increments and either repeats or reverses it. This
mechanism generates diffusive, critical and superdiffusive regimes, depending on
the memory parameter. Subsequent mathematical work has developed central limit
theorems, martingale methods, urn connections and variants with restricted
memory, random step sizes and delays \citep{colette2017central,bercu2018martingale,baur2016elephant,gut2021variations,gut2023review,dedecker2023rates,aguech2024clt,aguech2024delays}.

Artificial intelligence (AI), particularly online learning and reinforcement
learning, is also concerned with sequential decisions under uncertainty. In many
algorithms, exploration is controlled by a parameter that is updated from the
past. Yet the usual random-walk component of such methods is often Markovian or
externally scheduled. This paper studies a complementary question: what happens
when the memory parameter of an ERW is itself produced by a learning rule?

The contribution is a learning-controlled elephant random walk (LC-ERW). The
walker has complete elephant memory, but the probability of repeating a recalled
increment is chosen from the past by an arbitrary predictable learning
algorithm. This includes deterministic schedules, online estimators, stochastic
approximation recursions and bounded policy maps. The model is deliberately
minimal: it is rich enough to represent AI-style adaptation, but constrained
enough to support rigorous proofs.

The results are as follows. First, for any predictable memory policy, the
LC-ERW admits an exact martingale normalisation. Secondly, if the learned memory
coefficient converges fast enough to a diffusive value \(a<1/2\), then the
normalised position \(S_n/\sqrt{n}\) converges to a centred normal distribution
with variance \((1-2a)^{-1}\). Thirdly, a corollary gives a direct stability
condition for online-learning policies: any convergent policy satisfying a mild
summability condition inherits the diffusive ERW limit theorem. The numerical
section supplies reproducible R code that checks the predicted variance.

The manuscript is not a claim of peer-reviewed novelty. Rather, it states a
precise candidate contribution, proves it under transparent assumptions and
provides scripts that make the quantitative claims testable.

\section{Related Literature}

\citet{schutz2004elephants} introduced the ERW as a solvable non-Markovian
model with long-range memory. Limit theorems and related asymptotic properties
were developed by \citet{colette2017central}, while \citet{bercu2018martingale}
showed that martingale methods give a particularly effective route to strong
laws, quadratic strong laws, laws of the iterated logarithm and non-Gaussian
behaviour in the superdiffusive regime.

The connection with urn models was clarified by \citet{baur2016elephant}, and
extensions have since become an active topic. \citet{gut2021variations} and
\citet{gut2023review} discuss variants with modified memory mechanisms.
\citet{dedecker2023rates} study random step sizes and rates of convergence,
while \citet{aguech2024clt} and \citet{aguech2024delays} examine gradually
increasing memory, random step sizes and delays. These papers motivate the
present study but do not, to the best of our knowledge, formulate the memory
parameter itself as the output of a general predictable learning rule.

Table~\ref{tab:novelty} summarises the distinction. The proposed LC-ERW does
not claim to dominate these models; rather, it isolates a different source of
adaptivity. The increments remain the classical elephant increments, but the
repeat--reverse coefficient is updated from the observed past.

\begin{table}[h]
\centering
\caption{Position of the LC-ERW relative to common ERW variants.}
\label{tab:novelty}
\small
\begin{tabular}{p{0.25\textwidth}p{0.27\textwidth}p{0.36\textwidth}}
\toprule
Model family & Adapted component & Typical mathematical object \\
\midrule
Classical ERW & fixed memory parameter & non-Markovian walk \\
Restricted-memory ERW & recalled time set & memory window or memory growth \\
Random-step ERW & jump magnitude & random increments with ERW signs \\
Delayed ERW & move-or-stay decision & ERW with zeros \\
Pólya-type urn view & colour replacement dynamics & balanced urn representation \\
LC-ERW & predictable repeat probability & online policy \(p_n=\ell(\theta_n)\) \\
\bottomrule
\end{tabular}
\end{table}

\section{The Learning-Controlled Model}

Let \((\Omega,\F,\Pbb)\) be a probability space with filtration
\((\F_n)_{n\geq 1}\). The position of the walker after \(n\) steps is
\[
  S_n=\sum_{j=1}^n X_j,
  \qquad X_j\in\{-1,1\}.
\]
The first increment satisfies
\(\Pbb(X_1=1)=q\) and \(\Pbb(X_1=-1)=1-q\), where \(q\in[0,1]\).

For \(n\geq 1\), let \(K_n\) be conditionally uniform on
\(\{1,\ldots,n\}\) given \(\F_n\), and let \(p_n\) be an
\(\F_n\)-measurable random variable taking values in \((0,1)\). Conditional on
\(\F_n\) and \(K_n\), define
\[
  X_{n+1}=
  \begin{cases}
  X_{K_n}, & \text{with probability } p_n,\\
  -X_{K_n}, & \text{with probability } 1-p_n.
  \end{cases}
\]
The process \(p_n\) is the learning-controlled memory policy. In AI language,
it may be written as
\[
  p_n=\ell(\theta_n),
\]
where \(\theta_n\) is the state of an online learning algorithm and
\(\ell:\R^d\to(0,1)\) is a bounded policy link, for example a logistic map.
Set
\[
  a_n=2p_n-1.
\]
The coefficient \(a_n\in(-1,1)\) is the learned memory coefficient. Positive
values favour repetition of the recalled step, negative values favour reversal,
and \(a_n=0\) corresponds to no directional memory in conditional mean.

\begin{assumption}[Stable diffusive learning]\label{ass:stable}
There exist constants \(a\in(-1,1/2)\), \(\varepsilon\in(0,1)\),
\(\bar a<1/2\) and \(D<\infty\) such that, almost surely,
\[
  a_n\in[-1+\varepsilon,\,1-\varepsilon]
  \quad\text{and}\quad a_n\leq\bar a
  \quad\text{for all } n,
\]
and
\[
  \sum_{n=1}^{\infty}\frac{|a_n-a|}{n}\leq D.
\]
\end{assumption}

\begin{remark}
The summability condition is satisfied whenever \(a_n\to a\) at any polynomial
rate, for instance \(|a_n-a|=O(n^{-\rho})\) with \(\rho>0\). It is deliberately
stated pathwise, because the learning rule may be deterministic, data-driven or
generated by a stochastic approximation algorithm.
\end{remark}

\section{Main Results}

\begin{theorem}[Exact martingale normalisation]\label{thm:martingale}
Define \(A_1=1\) and, for \(n\geq 1\),
\[
  A_{n+1}=\prod_{k=1}^{n}\left(1+\frac{a_k}{k}\right).
\]
Then
\[
  M_n=\frac{S_n}{A_n},\qquad n\geq 1,
\]
is a martingale with respect to \((\F_n)\).
\end{theorem}

\begin{proof}
Since \(K_n\) is conditionally uniform on \(\{1,\ldots,n\}\),
\[
  \E[X_{K_n}\mid \F_n]=\frac{1}{n}\sum_{j=1}^{n}X_j=\frac{S_n}{n}.
\]
Using the conditional repeat--reverse rule,
\[
  \E[X_{n+1}\mid \F_n]
  =(p_n-(1-p_n))\frac{S_n}{n}
  =a_n\frac{S_n}{n}.
\]
Therefore
\[
  \E[S_{n+1}\mid \F_n]
  =S_n+a_n\frac{S_n}{n}
  =\left(1+\frac{a_n}{n}\right)S_n.
\]
By definition, \(A_{n+1}=(1+a_n/n)A_n\). Hence
\[
  \E\left[\frac{S_{n+1}}{A_{n+1}}\middle|\F_n\right]
  =\frac{S_n}{A_n}=M_n.
\]
The increments are bounded after normalisation on every finite horizon, so
integrability is immediate.
\end{proof}

\begin{lemma}[Asymptotics of the normalising product]\label{lem:product}
Under Assumption~\ref{ass:stable}, there is an almost surely finite random
variable \(L\in(0,\infty)\) such that
\[
  A_n=L n^a(1+o(1))\qquad\text{almost surely.}
\]
Moreover, there exist deterministic constants \(0<c_A<C_A<\infty\) such that
\[
  c_A n^a\leq A_n\leq C_A n^a
  \qquad\text{for all } n\geq1
\]
on the event where Assumption~\ref{ass:stable} holds.
\end{lemma}

\begin{proof}
Write
\[
  \log A_n-a\log n
  =
  \sum_{k=1}^{n-1}\left\{
  \log\left(1+\frac{a_k}{k}\right)-\frac{a}{k}
  \right\}
  +a\left(\sum_{k=1}^{n-1}\frac1k-\log n\right).
\]
The last term converges because the harmonic numbers satisfy
\(\sum_{k=1}^{n-1}k^{-1}-\log n\to\gamma\), where \(\gamma\) is Euler's
constant. For the first term, decompose
\[
  \log\left(1+\frac{a_k}{k}\right)-\frac{a}{k}
  =
  \frac{a_k-a}{k}
  +r_k,
  \qquad
  r_k=\log\left(1+\frac{a_k}{k}\right)-\frac{a_k}{k}.
\]
Because \(a_k\in[-1+\varepsilon,1-\varepsilon]\), the Taylor remainder is
uniformly bounded: there is \(C_\varepsilon<\infty\) such that
\(|r_k|\leq C_\varepsilon k^{-2}\) for all \(k\). Hence
\(\sum_k r_k\) converges absolutely, while
\(\sum_k (a_k-a)/k\) converges absolutely by Assumption~\ref{ass:stable}. Thus
\(\log A_n-a\log n\) converges to a finite real limit, say \(\log L\), which
proves \(A_n/(n^a)\to L\in(0,\infty)\).

The same estimates give deterministic bounds. Indeed,
\[
  \left|\log A_n-a\log n\right|
  \leq
  D+C_\varepsilon\sum_{k=1}^\infty k^{-2}
  +|a|\sup_{m\geq1}\left|\sum_{k=1}^{m}\frac1k-\log(m+1)\right|.
\]
Exponentiating this bound gives the constants \(c_A\) and \(C_A\).
\end{proof}

\begin{lemma}[Second moment control]\label{lem:second}
Under Assumption~\ref{ass:stable}, there is a deterministic constant \(C_2\)
such that
\[
  \sup_{n\geq 1}\frac{\E[S_n^2]}{n}\leq C_2.
\]
\end{lemma}

\begin{proof}
Let \(V_n=\E[S_n^2]\). Since \(X_{n+1}^2=1\),
\[
  \E[S_{n+1}^2\mid\F_n]
  =S_n^2+2S_n\E[X_{n+1}\mid\F_n]+1
  =\left(1+\frac{2a_n}{n}\right)S_n^2+1.
\]
Taking expectations and using \(a_n\leq\bar a<1/2\) gives
\[
  V_{n+1}\leq \left(1+\frac{2\bar a}{n}\right)V_n+1.
\]
Set \(\beta=2\bar a<1\). Iterating the recursion yields
\[
  V_n
  \leq
  V_1\prod_{j=1}^{n-1}\left(1+\frac{\beta}{j}\right)
  +\sum_{k=1}^{n-1}
  \prod_{j=k+1}^{n-1}\left(1+\frac{\beta}{j}\right).
\]
There is \(C_\beta<\infty\) such that
\(\prod_{j=k+1}^{n-1}(1+\beta/j)\leq C_\beta(n/k)^\beta\). Therefore
\[
  V_n\leq C_\beta n^\beta
  +C_\beta n^\beta\sum_{k=1}^{n-1}k^{-\beta}
  \leq C_2 n,
\]
because \(\beta<1\). This proves the claim.
\end{proof}

\begin{lemma}[Predictable quadratic variation]\label{lem:bracket}
Under Assumption~\ref{ass:stable},
\[
  \frac{A_n^2}{n}\langle M\rangle_n
  \xrightarrow{\Pbb}
  \frac{1}{1-2a}.
\]
\end{lemma}

\begin{proof}
The martingale increments are
\[
  \Delta M_{k+1}
  =
  \frac{X_{k+1}-a_kS_k/k}{A_{k+1}}.
\]
Consequently,
\[
  \langle M\rangle_n
  =
  \sum_{k=1}^{n-1}
  \frac{\E[(X_{k+1}-a_kS_k/k)^2\mid\F_k]}{A_{k+1}^2}
  =
  \sum_{k=1}^{n-1}
  \frac{1-a_k^2S_k^2/k^2}{A_{k+1}^2}.
\]
We first evaluate the leading term. Put \(B_n=A_n/n^a\). By
Lemma~\ref{lem:product}, \(B_n\to L\in(0,\infty)\) almost surely. Hence
\[
  \frac{A_n^2}{n}\sum_{k=1}^{n-1}\frac{1}{A_{k+1}^2}
  =
  B_n^2 n^{2a-1}
  \sum_{k=1}^{n-1}\frac{1}{B_{k+1}^2(k+1)^{2a}}.
\]
Since \(B_{k+1}^{-2}\to L^{-2}\) and \(2a<1\), Cesaro summation for regularly
varying sequences gives
\[
  n^{2a-1}
  \sum_{k=1}^{n-1}\frac{1}{B_{k+1}^2(k+1)^{2a}}
  \longrightarrow
  \frac{L^{-2}}{1-2a}.
\]
Multiplication by \(B_n^2\to L^2\) gives
\[
  \frac{A_n^2}{n}\sum_{k=1}^{n-1}\frac{1}{A_{k+1}^2}
  \longrightarrow
  \frac{1}{1-2a}
  \qquad\text{almost surely.}
\]

It remains to show that the correction term is negligible. By
Lemma~\ref{lem:product}, \(|a_k|\leq1\) and Lemma~\ref{lem:second},
\[
\begin{aligned}
  \E\left[
  \frac{A_n^2}{n}\sum_{k=1}^{n-1}
  \frac{a_k^2S_k^2/k^2}{A_{k+1}^2}
  \right]
  &\leq
  C n^{2a-1}
  \sum_{k=1}^{n-1} k^{-2a-2}\E[S_k^2]  \\
  &\leq
  C n^{2a-1}
  \sum_{k=1}^{n-1} k^{-2a-1}.
\end{aligned}
\]
The final expression tends to zero in all cases \(a<1/2\): if \(a>0\), the
sum is bounded; if \(a=0\), it is \(O(\log n)\); and if \(a<0\), it is
\(O(n^{-2a})\), giving \(O(n^{-1})\). Thus the correction term converges to
zero in \(L^1\), and hence in probability. The lemma follows.
\end{proof}

\begin{lemma}[Lindeberg condition]\label{lem:lindeberg}
Under Assumption~\ref{ass:stable}, for every \(\delta>0\),
\[
  \frac{1}{\langle M\rangle_n}
  \sum_{k=1}^{n-1}
  \E\left[(\Delta M_{k+1})^2
  \1_{\{|\Delta M_{k+1}|>\delta\langle M\rangle_n^{1/2}\}}
  \middle|\F_k\right]
  \xrightarrow{\Pbb}0.
\]
\end{lemma}

\begin{proof}
Because \(|X_{k+1}|\leq1\), \(|a_k|\leq1\), and \(|S_k|\leq k\),
\[
  |\Delta M_{k+1}|
  =
  \frac{|X_{k+1}-a_kS_k/k|}{A_{k+1}}
  \leq \frac{2}{A_{k+1}}.
\]
Lemma~\ref{lem:product} gives \(A_k\asymp k^a\). If \(a\geq0\), then
\(\max_{k\leq n}|\Delta M_k|=O(1)\); if \(a<0\), then
\(\max_{k\leq n}|\Delta M_k|=O(n^{-a})\). On the other hand,
Lemma~\ref{lem:bracket} implies
\[
  \langle M\rangle_n^{1/2}
  \asymp n^{1/2-a}
  \qquad\text{in probability.}
\]
Therefore
\[
  \frac{\max_{k\leq n}|\Delta M_k|}{\langle M\rangle_n^{1/2}}
  \xrightarrow{\Pbb}0.
\]
With probability tending to one, every indicator in the Lindeberg sum is then
zero, which proves the assertion.
\end{proof}

\begin{theorem}[Diffusive central limit theorem]\label{thm:clt}
Under Assumption~\ref{ass:stable},
\[
  \frac{S_n}{\sqrt{n}}
  \ \xrightarrow{\ \mathcal{D}\ }\ 
  \mathcal{N}\left(0,\frac{1}{1-2a}\right).
\]
\end{theorem}

\begin{proof}
The martingale \(M_n=S_n/A_n\) has increments satisfying the conditional
Lindeberg condition by Lemma~\ref{lem:lindeberg}. Since
\(\langle M\rangle_n\to\infty\) in probability by Lemmas~\ref{lem:product} and
\ref{lem:bracket}, the martingale central limit theorem
\citep[Theorem 3.2, for instance]{hall1980martingale} gives
\[
  \frac{M_n-M_1}{\langle M\rangle_n^{1/2}}
  \xrightarrow{\mathcal D}\mathcal N(0,1).
\]
Moreover \(M_1/\langle M\rangle_n^{1/2}\to0\) in probability. Hence
\[
  \frac{M_n}{\langle M\rangle_n^{1/2}}
  \xrightarrow{\mathcal D}\mathcal N(0,1).
\]
Lemma~\ref{lem:bracket} states that
\[
  \frac{A_n^2}{n}\langle M\rangle_n
  \xrightarrow{\Pbb}
  \frac{1}{1-2a}.
\]
Since \(S_n=A_nM_n\), Slutsky's theorem yields
\[
  \frac{S_n}{\sqrt n}
  =
  \frac{M_n}{\langle M\rangle_n^{1/2}}
  \left(\frac{A_n^2\langle M\rangle_n}{n}\right)^{1/2}
  \xrightarrow{\mathcal D}
  \mathcal N\left(0,\frac{1}{1-2a}\right).
\]
\end{proof}

\begin{corollary}[Stability of convergent learning policies]\label{cor:ai}
Let \(\theta_n\in\R^d\) be an \(\F_n\)-predictable learning state and let
\(\ell:\R^d\to(\varepsilon,1-\varepsilon)\) be Lipschitz. Put
\[
  p_n=\ell(\theta_n),\qquad a_n=2p_n-1.
\]
If \(\theta_n\to\theta_\star\) almost surely and there is
\(D_\theta<\infty\) such that
\[
  \sum_{n=1}^\infty \frac{\|\theta_n-\theta_\star\|}{n}\leq D_\theta
  \qquad\text{almost surely},
\]
with \(a=2\ell(\theta_\star)-1<1/2\), and if there exists
\(\bar a<1/2\) such that \(2\ell(\theta_n)-1\leq\bar a\) for all \(n\), then
the LC-ERW satisfies the central limit theorem in Theorem~\ref{thm:clt}.
\end{corollary}

\begin{proof}
By Lipschitz continuity,
\[
  |a_n-a|\leq 2\,\operatorname{Lip}(\ell)\|\theta_n-\theta_\star\|.
\]
The displayed summability condition implies Assumption~\ref{ass:stable}, and
Theorem~\ref{thm:clt} applies.
\end{proof}

\begin{proposition}[A damped empirical-drift learner]\label{prop:damped}
Fix \(\theta_\star\in\R\), \(c,\beta\in\R\), \(\rho>0\), \(\delta>0\), and
let \(\ell:\R\to(\varepsilon,1-\varepsilon)\) be Lipschitz. Define the
predictable learning state
\[
  \theta_n
  =
  \theta_\star+c(n+1)^{-\rho}
  +\beta\frac{S_n}{n^{1+\delta}},
  \qquad n\geq1,
\]
and set \(p_n=\ell(\theta_n)\). Suppose that \(a=2\ell(\theta_\star)-1<1/2\)
and that \(2\ell(\theta_n)-1\leq\bar a<1/2\) for all \(n\). Then the resulting
LC-ERW satisfies
\[
  \frac{S_n}{\sqrt n}
  \xrightarrow{\mathcal D}
  \mathcal N\left(0,\frac{1}{1-2a}\right).
\]
\end{proposition}

\begin{proof}
The rule is predictable because \(S_n\) is \(\F_n\)-measurable. Since
\(|S_n|\leq n\),
\[
  |\theta_n-\theta_\star|
  \leq |c|(n+1)^{-\rho}+|\beta|n^{-\delta}.
\]
Therefore
\[
  \sum_{n=1}^{\infty}\frac{|\theta_n-\theta_\star|}{n}
  \leq
  |c|\sum_{n=1}^{\infty}\frac{1}{n(n+1)^\rho}
  +|\beta|\sum_{n=1}^{\infty}\frac{1}{n^{1+\delta}}
  <\infty.
\]
The result follows from Corollary~\ref{cor:ai}. This proposition gives a
concrete learning rule whose correction term uses the observed empirical drift
but is damped strongly enough that the learned memory coefficient remains in
the stable diffusive class.
\end{proof}

\section{Examples of Learning Rules}

\paragraph{Deterministic annealing.}
Let \(p_n=p_\star+c(n+1)^{-\rho}\), truncated to
\((\varepsilon,1-\varepsilon)\), where \(p_\star<3/4\) and \(\rho>0\).
Then \(a_n\to a=2p_\star-1<1/2\) and the summability condition holds.

\paragraph{Online logistic policy.}
Let \(p_n=\ell(\theta_n)\), where
\(\ell(x)=\varepsilon+(1-2\varepsilon)/(1+\exp(-x))\), and suppose
\(\theta_n\) is an online estimator converging to \(\theta_\star\) at a
polynomial rate. The corollary gives the limiting law whenever
\(\ell(\theta_\star)<3/4\).

\paragraph{Damped empirical-drift learning.}
The rule in Proposition~\ref{prop:damped} is deliberately simple: the learner
uses the current empirical drift \(S_n/n\), but the term is multiplied by
\(n^{-\delta}\). This keeps the policy genuinely data-dependent while making
the stability condition verifiable without an additional stochastic
approximation theorem.

\paragraph{Stochastic approximation policies.}
More elaborate controllers can be analysed by first proving convergence of
\(\theta_n\) through standard stochastic approximation arguments
\citep{kushner2003stochastic}, and then applying Corollary~\ref{cor:ai}. This
two-step route is useful when the policy parameter is updated by a projected
Robbins--Monro recursion or by a bounded online estimator.

\section{Numerical Experiments}

The accompanying R scripts simulate the LC-ERW for several policies:

\begin{enumerate}
  \item a fixed-memory benchmark;
  \item a deterministic annealing policy;
  \item a logistic learning policy with polynomial convergence;
  \item the damped empirical-drift learner in Proposition~\ref{prop:damped}.
\end{enumerate}

For each policy, the scripts estimate the empirical variance of
\(S_n/\sqrt n\), compare it with \((1-2a)^{-1}\), and report Monte Carlo
standard errors, confidence intervals, skewness, excess kurtosis and
Kolmogorov--Smirnov diagnostics against the predicted Gaussian law. The scripts
also generate QQ-plot PDFs for visual inspection. These computations do not
prove the theorem; they check whether finite sample behaviour is consistent
with the asymptotic prediction and how quickly the variance ratio approaches
one as \(n\) grows.

\section{Discussion}

The LC-ERW shows that an AI-style adaptive memory policy can be incorporated
into an ERW without destroying its martingale backbone. The main theorem says
that, in the diffusive regime, stable learning is asymptotically equivalent to
freezing the memory parameter at its learned limit. This is useful in two
directions. For probability theory, it extends ERW analysis from fixed or
externally prescribed parameters to predictable data-driven policies. For AI,
it supplies a rare case in which an adaptive exploration mechanism with complete
memory has a transparent asymptotic law.

Several limitations remain. The theorem does not cover critical learning
limits \(a=1/2\), superdiffusive limits \(a>1/2\), multidimensional ERW, or
learning rules whose parameter does not converge. These cases are likely to
require different normalisations and may lead to non-Gaussian limits. Another
open direction is statistical inference: estimating the learned memory
trajectory \(p_n\) from observed LC-ERW paths.

\section{Conclusion}

This paper formulated a learning-controlled elephant random walk and proved
that predictable adaptive memory policies admit an exact martingale
normalisation. Under a stable diffusive learning condition, the process obeys a
central limit theorem with variance determined by the learned limiting memory
coefficient. The result gives a rigorous and reproducible starting point for
connecting elephant random walks with AI-controlled stochastic exploration.

\bibliographystyle{apalike}
\bibliography{references}

\end{document}
